\clearpage
\VerdierPhysicalPageAnchor{0156}%
\begingroup\small
\setlength{\parskip}{0pt}\setlength{\emergencystretch}{1.5em}

\VerdierSemanticAnchor{ch3-3-2-3-hom-map}
$(X,Y,f)$ と $(X',Y',f')$ を $\mathrm A(\mathcal D)$ の二対象とする。命題
(3.2.3) を示すには、関手 $\operatorname{Im}$ によって定まる準同型
\[
\operatorname{Hom}_{\mathrm A(\mathcal D)}((X,Y,f),(X',Y',f'))
\longrightarrow
\operatorname{Hom}_{\mathcal D^\wedge}(\operatorname{Im}(f),\operatorname{Im}(f'))
\tag{3.2.3.2}
\]
が同型であることを示せばよい。

\VerdierSemanticAnchor{ch3-3-2-3-injectivity}
まず単射性を示す。
$(g;g')\colon(X,Y,f)\to(X',Y',f')$ が
$\mathcal{F}\ell_1(\mathcal D)$ の射であり、
$\operatorname{Im}((g;g'))=0$ を満たすなら
\[
\operatorname{Im}((f';\operatorname{id}_{Y'}))\circ
\operatorname{Im}((g;g'))\circ
\operatorname{Im}((\operatorname{id}_X;f))=h(g'f)=0.
\]
$h$ は充満忠実なので $g'f=0$、ゆえに $(g;g')$ は無視可能である。
全射性は次の補題から従う。

\VerdierSemanticAnchor{ch3-3-2-4}
\medskip\noindent\textbf{補題 3.2.4.}\quad
\textit{$\alpha\colon\operatorname{Im}(f)\to\operatorname{Im}(f')$ を
$\mathcal D^\wedge$ の射とする。次の図式を可換にする $\mathcal D$ の二つの射
$g\colon X\to X'$ と $g'\colon Y\to Y'$ が存在する：}
\VerdierSemanticAnchor{ch3-3-2-4-1}
\[
\begin{tikzcd}[column sep=6em,row sep=2.3em]
h(X)\arrow[r,"\operatorname{Im}((\operatorname{id}_X;f))"]\arrow[d,"h(g)"']&
\operatorname{Im}(f)\arrow[d,"\alpha"]\\
h(X')\arrow[r,"\operatorname{Im}((\operatorname{id}_{X'};f'))"']&\operatorname{Im}(f')
\end{tikzcd}\tag{3.2.4.1}
\]
\VerdierSemanticAnchor{ch3-3-2-4-2}
\[
\begin{tikzcd}[column sep=5em,row sep=2.3em]
h(X)\arrow[r,"h(f)"]\arrow[d,"h(g)"']&h(Y)\arrow[d,"h(g')"]\\
h(X')\arrow[r,"h(f')"']&h(Y')
\end{tikzcd}\tag{3.2.4.2}
\]
\VerdierSemanticAnchor{ch3-3-2-4-consequence}
実際、(3.2.3.1) と $\operatorname{Im}((\operatorname{id}_X;f))$ が
$\mathcal D^\wedge$ におけるエピ射
であることから、補題は $\alpha=\operatorname{Im}((g;g'))$ を与える。
\endgroup
